\section{Conclusion}
\label{sec:conclusion}

For the complete total-spin-zero sector of an even finite spin-$1/2$ system, immunity to removal of one known particle is established prior art in the collective-symmetry DFS literature, notably Migda\l{} and Banaszek \cite{MigdalBanaszek2011}. Their retained-component orthogonality and equal-mixture construction also provide an equivalent form of the $W_i$ support factorization used here. We retain a self-contained channel formulation because it is convenient for the dynamical application; neither the loss-immunity mechanism nor the syndrome/logical factorization is claimed as new QEC/DFS mathematics.

The full-algebra recovery completion and the induced update $\Phi=\Lambda\circ\Ad_U\circ\R$ are standard consequences of exact recoverability. In the reviewed finite singlet-preserving constructions, these facts yield exact fixed per-$N$ autonomous records at $N=8,10,12$, specifically the reviewed $J2,J4,J6$ realizations.

These records are \emph{reconstructive}: their exact autonomy arises because they retain all information required to reconstruct the microscopic code state. That distinction organizes the interpretation of the reconciled manuscript: the reviewed finite records realize the reconstructive case, while autonomous closure after genuine information loss remains a separate question. This use of the distinction is a BCRD synthesis/application rather than a claim for a new general channel classification.

Accordingly, the contribution boundary is explicit. Prior art and standard consequences supply one-particle-loss immunity, the equivalent retained-support factorization, full-algebra recovery completion, and recovery-induced record closure. The BCRD-specific contribution is the finite $J2/J4/J6$ realization together with its record-dynamics synthesis. The finite constructions do not establish strict minimality, minimum-support growth, a common cross-$N$ map, asymptotic closure, continuum dynamics, geometry, gravity, or quantum gravity.
